\chapter{Abbildungen}

\begin{itemize}
\item Abbildung \ref{abb:grenze}: Ganzes und Grenze. Topisch und chronisch Kontinuierliches; der Chiasmus der Schnitte (nach Brentano).
\item Abbildung \ref{abb:faser}: Faser und Nebeneinander. Lichtkegel, Weltlinie als Faser, Faserbündel als Raum, Gleichzeitigkeitsflächen für verschiedene Werte von $\varepsilon$ (nach Reichenbach).
\item Abbildung \ref{abb:doppelgaenger}: Das symmetrische Universum. Doppelgänger, von außen ununterscheidbar, von innen indexikalisch unterschieden (nach Koch).
\item Abbildung \ref{abb:parameterraum}: Parameterraum und Koordinatenraum. Zwei Punkte auf einer Geraden oder ein Punkt in der Ebene; das Nahwirkungsprinzip zeichnet die Dimensionszahl aus (nach Reichenbach).
\item Abbildung \ref{abb:lichtjahr}: Das Lichtjahr als Urform der Länge. Signal, Reflexion, Rückkehr; Länge aus halber Laufzeit (nach Reichenbach).
\item Abbildung \ref{abb:uebergang}: Der Übergang und seine Richtung. Begrifflicher Aufstieg Punkt--Linie--Fläche--Raum--Zeit (Hegel) und Erfahrungsweg Zeit--Jetzt--Ort--Gegend (Guzzoni).
\item Abbildung \ref{abb:weile}: Die Weile im Fluss. Weile als Fläche gegen die Folge der Jetztpunkte; Ort als Augenblick des Raumes (nach Guzzoni).
\end{itemize}
